%%% WARNING %%%
%%% This file was automatically generated by SAPlugin,
%%% and will be overwritten next time.

% macro for maximum width and height
\makeatletter\def\maxwidth#1{\ifdim\Gin@nat@width>#1 #1\else\Gin@nat@width\fi}\makeatother
\makeatletter\def\maxheight#1{\ifdim\Gin@nat@height>#1 #1\else\Gin@nat@height\fi}\makeatother
\newlength{\txtheight}
\setlength{\txtheight}{\textheight}

\tikzset{
    hyperlink node/.style={
        alias=sourcenode,
        append after command={
            let     \p1 = (sourcenode.north west),
                \p2=(sourcenode.south east),
                \n1={\x2-\x1},
                \n2={\y1-\y2} in
            node [inner sep=0pt, outer sep=0pt,anchor=north west,at=(\p1)] {\hyperref[#1]{\XeTeXLinkBox{\phantom{\rule{\n1}{\n2}}}}}
                    %xelatex needs \XeTeXLinkBox, won't create a link unless it
                    %finds text --- rules don't work without \XeTeXLinkBox.
                    %Still builds correctly with pdflatex and lualatex
        }
    }
}


%%%%%%%%%%%%%%%%%%%%%%%%%%%%%%%%%%%%%%%%%%%%%%%%%%%%
%%% CS
%%%%%%%%%%%%%%%%%%%%%%%%%%%%%%%%%%%%%%%%%%%%%%%%%%%%
\chapter{Client-Server View}
\minilof{}


%%% Context diagram for the client-server view (1108.0x196.0)

\begin{figure}[!htp]
  \centering
  \includegraphics[width=\maxwidth{\textwidth}]{ComponentDiagram-Context-diagram-for-the-client-server-view.pdf}
  \caption[Context diagram for the client-server view]{\label{diag:Component:Contextdiagramfortheclientserverview}Context diagram for the client-server view}
\end{figure}



%%% Primary diagram of the client-server view (2830.0x1242.0; rotated)
\begin{landscape}
\begin{figure}[!htp]
  \centering
  \includegraphics[width=\maxwidth{\textwidth}]{ComponentDiagram-Primary-diagram-of-the-client-server-view.pdf}
  \caption[Primary diagram of the client-server view]{\label{diag:Component:Primarydiagramoftheclientserverview}Primary diagram of the client-server view\\This architecture is not yet complete.
Functionality that still has to be worked out in more detail is captured by the \vpett{\nameref{comp:ComponentseHealthPlatformOtherFunctionality}} component in gray.}
\end{figure}
\end{landscape}


%%% Decomposition of the Av2CommunicationGateway (988.0x421.0)

\begin{figure}[!htp]
  \centering
  \includegraphics[width=\maxwidth{\textwidth}]{ComponentDiagram-Decomposition-of-the-Av2CommunicationGateway.pdf}
  \caption[Decomposition of the Av2CommunicationGateway]{\label{diag:Component:DecompositionoftheAv2CommunicationGateway}Decomposition of the Av2CommunicationGateway}
\end{figure}

%%% Decomposition of the Av2MonitoringService (1066.0x420.0)

\begin{figure}[!htp]
  \centering
  \includegraphics[width=\maxwidth{\textwidth}]{ComponentDiagram-Decomposition-of-the-Av2MonitoringService.pdf}
  \caption[Decomposition of the Av2MonitoringService]{\label{diag:Component:DecompositionoftheAv2MonitoringService}Decomposition of the Av2MonitoringService}
\end{figure}



%%%%%%%%%%%%%%%%%%%%%%%%%%%%%%%%%%%%%%%%%%%%%%%%%%%%
%%% DECOMPOSITION
%%%%%%%%%%%%%%%%%%%%%%%%%%%%%%%%%%%%%%%%%%%%%%%%%%%%
\chapter{Decomposition View}
\minilof{}


%%% Context diagram for the decomposition view (558.0x247.0)

\begin{figure}[!htp]
  \centering
  \includegraphics[width=\maxwidth{\textwidth}]{ComponentDiagram-Context-diagram-for-the-decomposition-view.pdf}
  \caption[Context diagram for the decomposition view]{\label{diag:Component:Contextdiagramforthedecompositionview}Context diagram for the decomposition view}
\end{figure}



%%% Decomposition of the Av2PatientGateway (1145.0x601.0)

\begin{figure}[!htp]
  \centering
  \includegraphics[width=\maxwidth{\textwidth}]{ComponentDiagram-Decomposition-of-the-Av2PatientGateway.pdf}
  \caption[Decomposition of the Av2PatientGateway]{\label{diag:Component:DecompositionoftheAv2PatientGateway}Decomposition of the Av2PatientGateway}
\end{figure}

%%% Decomposition of the MLModelManager (809.0x274.0)

\begin{figure}[!htp]
  \centering
  \includegraphics[width=\maxwidth{\textwidth}]{ComponentDiagram-Decomposition-of-the-MLModelManager.pdf}
  \caption[Decomposition of the MLModelManager]{\label{diag:Component:DecompositionoftheMLModelManager}Decomposition of the MLModelManager}
\end{figure}

%%% Decomposition of the P1NotificationDispatcher (743.0x574.0)

\begin{figure}[!htp]
  \centering
  \includegraphics[width=\maxwidth{\textwidth}]{ComponentDiagram-Decomposition-of-the-P1NotificationDispatcher.pdf}
  \caption[Decomposition of the P1NotificationDispatcher]{\label{diag:Component:DecompositionoftheP1NotificationDispatcher}Decomposition of the P1NotificationDispatcher}
\end{figure}

%%% Decomposition of the P1PatientDataService (912.0x575.0)

\begin{figure}[!htp]
  \centering
  \includegraphics[width=\maxwidth{\textwidth}]{ComponentDiagram-Decomposition-of-the-P1PatientDataService.pdf}
  \caption[Decomposition of the P1PatientDataService]{\label{diag:Component:DecompositionoftheP1PatientDataService}Decomposition of the P1PatientDataService}
\end{figure}

%%% Decomposition of the P1PhysicianCommandService (845.0x583.0)

\begin{figure}[!htp]
  \centering
  \includegraphics[width=\maxwidth{\textwidth}]{ComponentDiagram-Decomposition-of-the-P1PhysicianCommandService.pdf}
  \caption[Decomposition of the P1PhysicianCommandService]{\label{diag:Component:DecompositionoftheP1PhysicianCommandService}Decomposition of the P1PhysicianCommandService}
\end{figure}

%%% Decomposition of the RiskEstimationProcessor (712.0x341.0)

\begin{figure}[!htp]
  \centering
  \includegraphics[width=\maxwidth{\textwidth}]{ComponentDiagram-Decomposition-of-the-RiskEstimationProcessor.pdf}
  \caption[Decomposition of the RiskEstimationProcessor]{\label{diag:Component:DecompositionoftheRiskEstimationProcessor}Decomposition of the RiskEstimationProcessor}
\end{figure}



%%%%%%%%%%%%%%%%%%%%%%%%%%%%%%%%%%%%%%%%%%%%%%%%%%%%
%%% DEPLOYMENT
%%%%%%%%%%%%%%%%%%%%%%%%%%%%%%%%%%%%%%%%%%%%%%%%%%%%
\chapter{Deployment view}
\minilof{}


%%% Context diagram for the deployment view (1091.0x879.0)

\begin{figure}[!htp]
  \centering
  \includegraphics[width=\maxwidth{\textwidth}]{DeploymentDiagram-Context-diagram-for-the-deployment-view.pdf}
  \caption[Context diagram for the deployment view]{\label{diag:Deployment:Contextdiagramforthedeploymentview}Context diagram for the deployment view}
\end{figure}



%%% Primary deployment diagram (3310.0x1055.0; rotated)
\begin{landscape}
\begin{figure}[!htp]
  \centering
  \includegraphics[width=\maxwidth{\textwidth}]{DeploymentDiagram-Primary-deployment-diagram.pdf}
  \caption[Primary deployment diagram]{\label{diag:Deployment:Primarydeploymentdiagram}Primary deployment diagram\\This deployment is necessarily incomplete as the architecture is incomplete.
The functionality that is not worked out yet is captured by \vpett{\nameref{comp:ComponentseHealthPlatformOtherFunctionality}} component, currently deployed on TODONode.
This TODONode will also have to be replaced with the final nodes are required by the replacement of \vpett{\nameref{comp:ComponentseHealthPlatformOtherFunctionality}}.}
\end{figure}
\end{landscape}


%%% Pilot Deployment (200 patients, 5 clin models) (1528.0x1013.2827; rotated)
\begin{landscape}
\begin{figure}[!htp]
  \centering
  \includegraphics[width=\maxwidth{\textwidth}]{DeploymentDiagram-Pilot-Deployment-200-patients-5-clin-models-.pdf}
  \caption[Pilot Deployment (200 patients, 5 clin models)]{\label{diag:Deployment:PilotDeployment200patients5clinmodels}Pilot Deployment (200 patients, 5 clin models)\\This diagram shows a concrete initial pilot deployment for 200 patients with an average of 5 clinical models per patient.
This diagram is a concrete instance of a deployment for the specified number of patients and clinical models, therefore it also includes information from the context diagram.
Note that it is not correct to include this type of information in your primary deployment diagram.}
\end{figure}
\end{landscape}
%%% Development Test Deployment (20 patients, 3 clin models) (1110.0x830.0)

\begin{figure}[!htp]
  \centering
  \includegraphics[width=\maxwidth{\textwidth}]{DeploymentDiagram-Development-Test-Deployment-20-patients-3-clin-models-.pdf}
  \caption[Development Test Deployment (20 patients, 3 clin models)]{\label{diag:Deployment:DevelopmentTestDeployment20patients3clinmodels}Development Test Deployment (20 patients, 3 clin models)\\This diagram shows a simplified deployment for local development.
This diagram is a concrete instance of a deployment for the specified number of patients and clinical models, therefore it also includes information from the context diagram.
Note that it is not correct to include this type of information in your primary deployment diagram.}
\end{figure}



%%%%%%%%%%%%%%%%%%%%%%%%%%%%%%%%%%%%%%%%%%%%%%%%%%%%
%%% SEQUENCE
%%%%%%%%%%%%%%%%%%%%%%%%%%%%%%%%%%%%%%%%%%%%%%%%%%%%
\chapter{Process View}
\minilof{}


%%% Risk estimation process (1168.0x599.0)

\begin{figure}[!htp]
  \centering
  \begin{tikzpicture}
    \node[anchor=south west,inner sep=0] (image) at (0,0) {\includegraphics[width=\maxwidth{\textwidth}]{InteractionDiagram-Risk-estimation-process.pdf}};
    \begin{scope}[x={(image.south east)},y={(image.north west)}]
        \node [fit={(0.8759,1.0000) (0.9940,0.9399)},inner sep=0,hyperlink node={comp:ComponentseHealthPlatformClinicalModelDB}] {};
\node [fit={(0.3219,1.0000) (0.4418,0.9399)},inner sep=0,hyperlink node={comp:ComponentseHealthPlatformClinicalJobCreator}] {};
\node [fit={(0.0000,1.0000) (0.1164,0.9399)},inner sep=0,hyperlink node={comp:ComponentseHealthPlatformOtherFunctionality}] {};
\node [fit={(0.5291,1.0000) (0.6498,0.9399)},inner sep=0,hyperlink node={comp:ComponentseHealthPlatformClinicalModelCache}] {};
\node [fit={(0.6704,1.0000) (0.8048,0.9399)},inner sep=0,hyperlink node={comp:ComponentseHealthPlatformRiskEstimationCombiner}] {};
\node [fit={(0.1721,1.0000) (0.3099,0.9399)},inner sep=0,hyperlink node={comp:ComponentseHealthPlatformRiskEstimationScheduler}] {};
    \end{scope}
    \end{tikzpicture}
  \caption[Risk estimation process]{\label{diag:Interaction:Riskestimationprocess}Risk estimation process\\The scheduling and execution of a risk estimation triggered by the arrival of new sensor data. Note that if one or more clinical models in the \vpett{\nameref{comp:ComponentseHealthPlatformClinicalModelCache}} have been invalidated, all relevant clinical models are renewed.}
\end{figure}

%%% Compute clinical model result (1365.0x693.0)

\begin{figure}[!htp]
  \centering
  \begin{tikzpicture}
    \node[anchor=south west,inner sep=0] (image) at (0,0) {\includegraphics[width=\maxwidth{\textwidth}]{InteractionDiagram-Compute-clinical-model-result.pdf}};
    \begin{scope}[x={(image.south east)},y={(image.north west)}]
        \node [fit={(0.1392,1.0000) (0.2564,0.9481)},inner sep=0,hyperlink node={mod:ComponentseHealthPlatformRiskEstimationProcessorJobProcessor}] {};
\node [fit={(0.5458,1.0000) (0.6593,0.9481)},inner sep=0,hyperlink node={mod:ComponentseHealthPlatformRiskEstimationProcessorSensorDataRiskAssessment}] {};
\node [fit={(0.3993,1.0000) (0.5201,0.9481)},inner sep=0,hyperlink node={mod:ComponentseHealthPlatformRiskEstimationProcessorMLModelRetrieval}] {};
\node [fit={(0.2747,1.0000) (0.3919,0.9481)},inner sep=0,hyperlink node={mod:ComponentseHealthPlatformRiskEstimationProcessorJobMLModelSelector}] {};\node [fit={(0.9187,1.0000) (0.9993,0.9481)},inner sep=0,hyperlink node={int:InterfacesKVStore}] {};
\node [fit={(0.7509,1.0000) (0.8205,0.9481)},inner sep=0,hyperlink node={int:InterfacesMLAPI}] {};
\node [fit={(0.0000,1.0000) (0.0916,0.9481)},inner sep=0,hyperlink node={int:InterfacesFetchJobs}] {};
\node [fit={(0.8278,1.0000) (0.9011,0.9481)},inner sep=0,hyperlink node={int:InterfacesResults}] {};
\node [fit={(0.6740,1.0000) (0.7443,0.9481)},inner sep=0,hyperlink node={int:InterfacesMLModelMgmt}] {};
    \end{scope}
    \end{tikzpicture}
  \caption[Compute clinical model result]{\label{diag:Interaction:Computeclinicalmodelresult}Compute clinical model result\\This sequence diagram shows the detailed application and computation of the clinical model. The first step involves the retrieval of the appropriate \ref{data:DatatypesMLModel} (specified by the \ref{data:DatatypesClinicalModel}). The second step involves the calculation of the prediction using the retrieved \ref{data:DatatypesMLModel}s.
Note that, as this diagram depicts the internal flow between modules within the \vpett{\nameref{comp:ComponentseHealthPlatformRiskEstimationProcessor}}, the gray lifelines depict its (external) interfaces.}
\end{figure}

%%% MLModel Synchronization (1276.0x521.0)

\begin{figure}[!htp]
  \centering
  \begin{tikzpicture}
    \node[anchor=south west,inner sep=0] (image) at (0,0) {\includegraphics[width=\maxwidth{\textwidth}]{InteractionDiagram-MLModel-Synchronization.pdf}};
    \begin{scope}[x={(image.south east)},y={(image.north west)}]
        \node [fit={(0.0980,1.0000) (0.2328,0.9309)},inner sep=0,hyperlink node={mod:ComponentseHealthPlatformMLModelManagerMLModelStorageManager}] {};
\node [fit={(0.4099,1.0000) (0.5086,0.9309)},inner sep=0,hyperlink node={mod:ComponentseHealthPlatformMLModelManagerMLModelRepository}] {};
\node [fit={(0.2555,1.0000) (0.3871,0.9309)},inner sep=0,hyperlink node={mod:ComponentseHealthPlatformMLModelManagerMLModelUpdateProcessor}] {};\node [fit={(0.0000,1.0000) (0.0862,0.9309)},inner sep=0,hyperlink node={int:InterfacesMLModelMgmt}] {};
\node [fit={(0.9146,1.0000) (0.9992,0.9309)},inner sep=0,hyperlink node={int:InterfacesMLaaSService}] {};
\node [fit={(0.7649,1.0000) (0.8229,0.9309)},inner sep=0,hyperlink node={int:InterfacesMLaaSAPI}] {};
    \end{scope}
    \end{tikzpicture}
  \caption[MLModel Synchronization]{\label{diag:Interaction:MLModelSynchronization}MLModel Synchronization\\As the system operations new \ref{data:DatatypesMLModel}s may become available for use by the \ref{data:DatatypesClinicalModel}s. This diagram illustrates the synchronization logic to fetch any unavailable \ref{data:DatatypesMLModel}s and to store a local copy of the \ref{data:DatatypesMLModel}s in the \vpett{\nameref{mod:ComponentseHealthPlatformMLModelManagerMLModelRepository}} for faster retrieval and use in the sensor data risk assessment.
Note that, as this diagram depicts the internal flow between modules within the \vpett{\nameref{comp:ComponentseHealthPlatformMLModelManager}}, the gray lifelines depict its (external) interfaces.}
\end{figure}

%%% Processing MLModel updates (1353.0x839.0)

\begin{figure}[!htp]
  \centering
  \begin{tikzpicture}
    \node[anchor=south west,inner sep=0] (image) at (0,0) {\includegraphics[width=\maxwidth{\textwidth}]{InteractionDiagram-Processing-MLModel-updates.pdf}};
    \begin{scope}[x={(image.south east)},y={(image.north west)}]
        \node [fit={(0.1123,1.0000) (0.2483,0.9571)},inner sep=0,hyperlink node={mod:ComponentseHealthPlatformMLModelManagerMLModelStorageManager}] {};
\node [fit={(0.6327,1.0000) (0.7214,0.9571)},inner sep=0,hyperlink node={mod:ComponentseHealthPlatformMLModelManagerModelUpdater}] {};
\node [fit={(0.4686,1.0000) (0.5928,0.9571)},inner sep=0,hyperlink node={mod:ComponentseHealthPlatformMLModelManagerMLModelUpdateProcessor}] {};
\node [fit={(0.3245,1.0000) (0.4287,0.9571)},inner sep=0,hyperlink node={mod:ComponentseHealthPlatformMLModelManagerMLModelRepository}] {};\node [fit={(0.8145,1.0000) (0.8854,0.9571)},inner sep=0,hyperlink node={int:InterfacesMLAPI}] {};
\node [fit={(0.7354,1.0000) (0.8004,0.9571)},inner sep=0,hyperlink node={int:InterfacesMLaaSAPI}] {};
\node [fit={(0.9039,1.0000) (0.9993,0.9571)},inner sep=0,hyperlink node={int:InterfacesMLaaSService}] {};
\node [fit={(0.0000,1.0000) (0.0916,0.9571)},inner sep=0,hyperlink node={int:InterfacesMLModelMgmt}] {};
    \end{scope}
    \end{tikzpicture}
  \caption[Processing MLModel updates]{\label{diag:Interaction:ProcessingMLModelupdates}Processing MLModel updates\\\ref{data:DatatypesMLModelUpdate}s are made available to improve the existing \ref{data:DatatypesMLModel}s. This diagram illustrates the retrieval of these \ref{data:DatatypesMLModelUpdate}s and their application on the local \ref{data:DatatypesMLModel}. As updates to these models influence the risk assessments, they need to be triggered for each individual \ref{data:DatatypesMLModel} that requires an update.
Note that, as this diagram depicts the internal flow between modules within the \vpett{\nameref{comp:ComponentseHealthPlatformMLModelManager}}, the gray lifelines depict its (external) interfaces.}
\end{figure}

%%% Incoming sensor data (734.0x242.0)

\begin{figure}[!htp]
  \centering
  \begin{tikzpicture}
    \node[anchor=south west,inner sep=0] (image) at (0,0) {\includegraphics[width=\maxwidth{\textwidth}]{InteractionDiagram-Incoming-sensor-data.pdf}};
    \begin{scope}[x={(image.south east)},y={(image.north west)}]
        \node [fit={(0.2221,0.5868) (0.6662,0.3802)},inner sep=0,hyperlink node={diag:Interaction:Riskestimationprocess}] {};\node [fit={(0.4728,1.0000) (0.7112,0.8512)},inner sep=0,hyperlink node={comp:ComponentseHealthPlatformRiskEstimationCombiner}] {};
\node [fit={(0.2439,1.0000) (0.4632,0.8512)},inner sep=0,hyperlink node={comp:ComponentseHealthPlatformRiskEstimationScheduler}] {};
\node [fit={(0.0000,1.0000) (0.1703,0.8512)},inner sep=0,hyperlink node={comp:ComponentseHealthPlatformOtherFunctionality}] {};
    \end{scope}
    \end{tikzpicture}
  \caption[Incoming sensor data]{\label{diag:Interaction:Incomingsensordata}Incoming sensor data\\This diagram shows a small part of the processing of incoming sensor data. It does not yet depict the arrival of the sensor data, storage in a database, or any other processing, as this functionality is not yet worked out.}
\end{figure}

%%% Combining ClinicalModelJob results (1007.0x247.0)

\begin{figure}[!htp]
  \centering
  \begin{tikzpicture}
    \node[anchor=south west,inner sep=0] (image) at (0,0) {\includegraphics[width=\maxwidth{\textwidth}]{InteractionDiagram-Combining-ClinicalModelJob-results.pdf}};
    \begin{scope}[x={(image.south east)},y={(image.north west)}]
        \node [fit={(0.0000,1.0000) (0.1629,0.8543)},inner sep=0,hyperlink node={comp:ComponentseHealthPlatformRiskEstimationProcessor}] {};
\node [fit={(0.6564,1.0000) (0.7915,0.8543)},inner sep=0,hyperlink node={comp:ComponentseHealthPlatformOtherFunctionality}] {};
\node [fit={(0.2632,1.0000) (0.4191,0.8543)},inner sep=0,hyperlink node={comp:ComponentseHealthPlatformRiskEstimationCombiner}] {};
    \end{scope}
    \end{tikzpicture}
  \caption[Combining ClinicalModelJob results]{\label{diag:Interaction:CombiningClinicalModelJobresults}Combining ClinicalModelJob results}
\end{figure}
